% Nouveau EGA V (cote 25): archive cote 25, site page 131 (PDF page 132).
% typescript. Draft compared: none.
% AI-transcribed by Claude Opus 5.5 (claude-opus-5-5, adaptive thinking, effort high), Message Batches API, batch ega25b (msgbatch_01A7n5YuXubUPNK1zsepi8RG), 2026-11-13; not yet checked by eye.
\scanpage{25/131}
\mspage{34}

par $\underline{\mathcal{O}}_{X_o}(D_o)$ sur D$_o$. Dans le cas b), \unsure{On} peut supposer D sans points
fixes, et il est plus commode d'interpréter la situation en termes de
morphismes dans P (cf 4.1.), de sorte que G devient le foncteur
\underline{Hom}$_S$(X,P), et \hand{$\mathcal{V}$}$_{w_o}$ sera l'espace \xout{\illegible{1}} Hom$_{\underline{O}_{X_o}}$$(f_o^{*}(\underline{\hand{$\Omega$}}^{1}_{P_o/k}), \underline{O}_{P_o})$.
Dans l'un et l'autre cas, on aura un homomorphisme naturel
\hand{$\mathcal{V}$}$_{u_o}\otimes_k V$ \hand{$\to$} \hand{$\mathcal{V}$}$_{w_o}\otimes_k V$ (où \hand{$\mathcal{V}$}$_{u_o}$ = dual de $u_o^{*}($\hand{$\Omega$}$^{1}_{F_o/k})$ ) exprimant le passage
de la déviation \hand{$\delta$} à la déviation \hand{$\delta$}' correspondante par application de h,
et l'injectivité de cette application résulte\xout{ra} encore de l'injectivité
de \hand{$\mathcal{V}$}$_{u_o}$\hand{$\to$} \hand{$\mathcal{V}$}$_{w_o}$ , qui elle provient\xout{ra} du fait que $h_o$: $F_o$\hand{$\to$}$G_o$ est un
monomorphisme.
\margin{$/$}

Pratiquement, il ne semble guère possible de rédiger cette dernière
partie de la démonstration sans référer aux petits calculs du par.2\hand{5}
(qu'il n'est pas question de refaire ici dans un cas particulier). On
signalera que cela ne donne pas lieu à des cercles vicieux, car le par.
2\hand{5} \xout{xxx} et les calculs qu'on y \xout{xxxxx} développe ne dépend que des sorites
de calcul différentiel du par.16, et d'ailleurs \hand{4}.3. ne sera plus
utilisé dans le Chap IV, sauf peut-être dans les deux n$^{o}$s qui suivent.

L'intérêt de \hand{4}.3.\inserted{a)} est de montrer que sous les conditions \deleted{de a)}\inserted{\unsure{initiales}},
le fibré projectif paramétrisant peut s'interpréter comme un sous-foncteur
de Div$_{X/S}$ , ce qui permet d'interpréter de façon intrinsèque la notion
de classe (à isomorphisme près sur le \xout{schéma} fibré projectif paramétrisa-%
teur) de système linéaire sur X/S, e\xout{n termes d'un} comme étant un sous-
foncteur $\overset{\hand{$\vee$}}{\text{P}}$ de Div$_{X/S}$ qui satisfait certaines propriétés (\unsure{savoir} est
représentable par un fibré projectif, et la famille de diviseurs définie
par l'injection c anonique de ce dernier dans Div$_{X/S}$ est linéaire au s ens
du n$^{o}$\over{4}{3}), ce qui est \xout{\illegible{1}}\inserted[typed]{essentiellement} le point de vue classique (où un\xout{e}
système linéaire de diviseurs est défini comme un \uline{ensemble} de diviseurs

\margin{131}
